% Paths of the Wild: player options in the style of the 2024 core books.
%
% Shows: SRD 5.2.1 attribution, text that changes with the SRD version
% (\DndIfSRD), modern stat blocks, the print background (footer scroll, no
% parchment), justified text, colored links, a theme color, a class table, a
% spell list and an index.
%
% Build from the template folder (pdfLaTeX, free fonts):
%   make all-editions BOOK=examples/player-options/player-options.tex
% and for the 2014 rules, from the same source:
%   make all-editions BOOK=examples/player-options/player-options.tex SRD=5.1
\documentclass[letterpaper,twoside,twocolumn,openany,nodeprecatedcode,
  srd=5.2.1,stats=modern,bg=print,justified,links=color,index]{dndbook}

\usepackage[english]{babel}
\usepackage[utf8]{inputenc}

\title{Paths of the Wild}
\author{A. Writer}
\DndSetMetadata{
  subject  = {Player options for rangers and druids},
  keywords = {D\&D, 5e, player options},
}

% The word the rules use for a character's people
\newcommand{\species}{\DndIfSRD{5.2}{species}{race}}

\begin{document}

\frontmatter

\DndCreditsPage[columns=1]{
  \DndCredit{design}{A. Writer}
  \DndCredit{development}{B. Developer}
  \DndCredit{layout}{The dndbook class}
  \DndCredit{Rules}{\DndIfSRD{5.2}{2024}{2014} rules, from SRD~\DndSRDVersion}
}

\tableofcontents

\mainmatter

\DndSetThemeColor[DmgCoral]

\chapter{The Warden}

Wardens\index{Warden} guard the wild places between settlements. Any \species{} can take up the path; most wardens learn it from an older warden\index{Warden!training} who has outlived their own patch of forest.

\begin{DndClassTable}[title=The Warden, header-rows=2]{ccXccc}
  \DndClassTableGroup{3}{3}{Spell Slots per Spell Level}
  \DndClassTableHeader{Level & Proficiency Bonus & Features & 1st & 2nd & 3rd}
  1st & +2 & Spellcasting, Wild Ward      & 2 & --- & --- \\
  2nd & +2 & Pathfinder                   & 2 & --- & --- \\
  3rd & +2 & Warden Oath                  & 3 & ---  & --- \\
  4th & +2 & Ability Score Improvement    & 3 & ---  & --- \\
  5th & +3 & Extra Attack                 & 4 & 2   & --- \\
  6th & +3 & Oath Feature                 & 4 & 2   & --- \\
  7th & +3 & Roots of the World           & 4 & 3   & --- \\
  8th & +3 & Ability Score Improvement    & 4 & 3   & --- \\
  9th & +4 & ---                          & 4 & 3   & 2   \\
\end{DndClassTable}

\section{Class Features}

\subsection{Wild Ward}
You can use a bonus action to mark a creature you can see within 60 feet of you. For the next minute, the marked creature has disadvantage on attack rolls against anyone but you.\index{Wild Ward}

\subsection{Pathfinder}
You and your companions cannot become lost except by magical means, and difficult terrain made of plants does not slow you.\index{Pathfinder}

\begin{DndComment}{Choosing a \DndIfSRD{5.2}{Species}{Race}}
  Wardens work best with a \species{} that has good Wisdom. \DndIfSRD{5.2}{In the 2024 rules your background, not your species, raises your ability scores.}{In the 2014 rules your race raises your ability scores.}
\end{DndComment}

\begin{DndSpellList}{Warden Spells}
  \DndSpellListLevel{0}{Druidcraft, Guidance, Thorn Whip}
  \DndSpellListLevel{1}{Entangle, Goodberry, Speak with Animals, Hunter's Mark}
  \DndSpellListLevel{2}{Barkskin, Pass without Trace, Spike Growth}
  \DndSpellListLevel{3}{Plant Growth, Wind Wall}
\end{DndSpellList}

\DndSpellHeader%
  {Warden's Bramble}
  {Level 1 Conjuration (Warden)}
  {Bonus action}
  {30 feet}
  {V, S}
  {Concentration, up to 1 minute}
Thorny vines burst from the ground in a 10-foot square you can see. The area is difficult terrain, and a creature that enters it for the first time on a turn takes \DndDice{1d6} piercing damage.\index{spells!Warden's Bramble}

\chapter{Companions}

A warden may bond with a beast of the wild, which fights at the warden's side.\index{companions}

\begin{DndMonster}{Ridge Elk}
  \DndMonsterType{Large Beast, Unaligned}

  \DndMonsterBasics[
      armor-class = {12},
      initiative  = {+2 (12)},
      hit-points  = {\DndDice{3d10 + 6}},
      speed       = {50 ft.},
    ]

  \DndMonsterAbilityScores[
      str = 16, dex = 14, con = 14,
      int = 2, wis = 12, cha = 5,
    ]

  \DndMonsterDetails[
      senses    = {Passive Perception 11},
      languages = {None},
      challenge = {1/2},
    ]

  \DndMonsterSection{Traits}
  \DndMonsterAction{Charge}
  If the elk moves at least 20 feet straight toward a target and then hits it with a Ram attack on the same turn, the target takes an extra \DndDice{2d6} bludgeoning damage.

  \DndMonsterSection{Actions}
  \DndMonsterMelee[
    name=Ram,
    mod=+5,
    dmg=\DndDice{1d6+3},
    dmg-type=bludgeoning,
  ]
\end{DndMonster}

Your companion acts on your initiative, right after you.\index{companions!initiative} See the \nameref{sec:bond} rules below.

\section{The Bond}\label{sec:bond}
If your companion drops to 0 hit points, you can spend a Hit Die and touch it to return it to 1 hit point.\index{companions!reviving}

\DndPrintIndex

\end{document}
